
\documentclass{article}
\usepackage{colortbl}
\usepackage{makecell}
\usepackage{multirow}
\usepackage{supertabular}

\begin{document}

\newcounter{utterance}

\centering \large Interaction Transcript for game `function\_detective', experiment `numbers\_passive\_labeled\_pairs', episode 33 with Qwen3.6{-}35B{-}A3B{-}FP8{-}without{-}reasoning.
\vspace{24pt}

{ \footnotesize  \setcounter{utterance}{1}
\setlength{\tabcolsep}{0pt}
\begin{supertabular}{c@{$\;$}|p{.15\linewidth}@{}p{.15\linewidth}p{.15\linewidth}p{.15\linewidth}p{.15\linewidth}p{.15\linewidth}}
    \# & \multicolumn{2}{c}{Player} && \multicolumn{2}{c}{Game Master} \\
    \hline

    \theutterance \stepcounter{utterance}  
    & & & \multicolumn{4}{p{0.6\linewidth}}{
        \cellcolor[rgb]{0.9,0.9,0.9}{
            \makecell[{{p{\linewidth}}}]{
                \texttt{\tiny{[P1$\langle$GM]}}
                \texttt{I have a hidden Python function for you to figure out.} \\
\\ 
\texttt{Your job is to infer what it does and then submit a correct Python implementation.} \\
\\ 
\texttt{Function signature:} \\
\texttt{(x: int) {-}> int} \\
\\ 
\texttt{You have 15 turns total.} \\
\\ 
\texttt{I reveal candidate input{-}output pairs one at a time and tell you whether each one is correct.} \\
\\ 
\texttt{You must make your best SOLVE attempt no later than the second{-}to{-}last turn.} \\
\texttt{Do not wait until the very end.} \\
\\ 
\texttt{You can use these actions:} \\
\texttt{{-} JSON format described at the end of the prompt} \\
\texttt{{-} SOLVE: ```python ... ```} \\
\\ 
\texttt{Do not add import statements in your guess.} \\
\texttt{Anything you might need is already available, so just write the function body in hypothesis and return the answer.} \\
\\ 
\texttt{Preloaded examples:} \\
\texttt{1. ((1000, 4), True)} \\
\\ 
\texttt{Scoring:} \\
\texttt{{-} correctness is required for any non{-}zero score} \\
\texttt{{-} in interactive modes, solving with fewer information requests is better} \\
\\ 
\texttt{Input formatting reminders:} \\
\texttt{{-} strings must be quoted} \\
\texttt{{-} booleans must be True/False} \\
\texttt{{-} use commas between multiple values} \\
\texttt{{-} do not write any extra action tags in the same message} \\
\\ 
\texttt{Format example:} \\
\texttt{You write:} \\
\texttt{\{'current\_mode': <probing/confirming/don't know; choose whichever best describes your current state in the task>, 'input\_rationale': <brief description of why you chose this hypothesis>, 'hypothesis': <either the exact string "unknown" if you have no guess yet OR your current best guess as runnable Python matching the given signature, with a one{-}line docstring stating what it does. Never write placeholder code, if you cannot commit to a guess, use "unknown". Match the given signature.} \\
\texttt{e.g.} \\
\texttt{def f(x):} \\
\texttt{	"""Returns the absolute value of x."""} \\
\texttt{	return x if x >= 0 else {-}x} \\
\\ 
\texttt{Use "unknown" if you have no guess.>, 'confidence\_score': <score between 1{-}100. 100 means highest confidence in your hypothesis>\}} \\
\\ 
\texttt{Output only one valid json per turn.} \\
            }
        }
    }
     \\ \\

    \theutterance \stepcounter{utterance}  
    & \multicolumn{4}{p{0.6\linewidth}}{
        \cellcolor[rgb]{0.561,0.737,0.561}{
            \makecell[{{p{\linewidth}}}]{
                \texttt{\tiny{[P1$\rangle$GM]}}
                \texttt{\{'current\_mode': 'probing', 'input\_rationale': 'I have only one data point: f(1000) = 4. This is not enough to determine the function. I need more examples to understand the pattern. I will try a few small integers to see if there is a simple arithmetic or digit{-}based relationship.', 'hypothesis': 'unknown', 'confidence\_score': 10\}} \\
            }
        }
    }
    & & \\ \\

    \theutterance \stepcounter{utterance}  
    & & & \multicolumn{4}{p{0.6\linewidth}}{
        \cellcolor[rgb]{0.9,0.9,0.9}{
            \makecell[{{p{\linewidth}}}]{
                \texttt{\tiny{[P1$\langle$GM]}}
                \texttt{Game Violation: ParseError: Response is not valid JSON.} \\
\texttt{Output ONLY:} \\
\texttt{One JSON with confidence\_score, current\_mode, input\_rationale, and hypothesis keys ...} \\
\texttt{OR} \\
\texttt{SOLVE: ```python ... ```} \\
            }
        }
    }
     \\ \\

    \theutterance \stepcounter{utterance}  
    & \multicolumn{4}{p{0.6\linewidth}}{
        \cellcolor[rgb]{0.561,0.737,0.561}{
            \makecell[{{p{\linewidth}}}]{
                \texttt{\tiny{[P1$\rangle$GM]}}
                \texttt{\{'current\_mode': 'probing', 'input\_rationale': 'I have only one data point: f(1000) = 4. This is not enough to determine the function. I need more examples to understand the pattern. I will try a few small integers to see if there is a simple arithmetic or digit{-}based relationship.', 'hypothesis': 'unknown', 'confidence\_score': 10\}} \\
            }
        }
    }
    & & \\ \\

    \theutterance \stepcounter{utterance}  
    & & & \multicolumn{4}{p{0.6\linewidth}}{
        \cellcolor[rgb]{0.9,0.9,0.9}{
            \makecell[{{p{\linewidth}}}]{
                \texttt{\tiny{[P1$\langle$GM]}}
                \texttt{Game Violation: ParseError: Response is not valid JSON.} \\
\texttt{Output ONLY:} \\
\texttt{One JSON with confidence\_score, current\_mode, input\_rationale, and hypothesis keys ...} \\
\texttt{OR} \\
\texttt{SOLVE: ```python ... ```} \\
            }
        }
    }
     \\ \\

    \theutterance \stepcounter{utterance}  
    & \multicolumn{4}{p{0.6\linewidth}}{
        \cellcolor[rgb]{0.561,0.737,0.561}{
            \makecell[{{p{\linewidth}}}]{
                \texttt{\tiny{[P1$\rangle$GM]}}
                \texttt{\{"current\_mode": "probing", "input\_rationale": "I have only one data point: f(1000) = 4. This is not enough to determine the function. I need more examples to understand the pattern. I will try a few small integers to see if there is a simple arithmetic or digit{-}based relationship.", "hypothesis": "unknown", "confidence\_score": 10\}} \\
            }
        }
    }
    & & \\ \\

    \theutterance \stepcounter{utterance}  
    & & & \multicolumn{4}{p{0.6\linewidth}}{
        \cellcolor[rgb]{0.9,0.9,0.9}{
            \makecell[{{p{\linewidth}}}]{
                \texttt{\tiny{[P1$\langle$GM]}}
                \texttt{OUTPUT: (({-}2, {-}1), False)} \\
            }
        }
    }
     \\ \\

    \theutterance \stepcounter{utterance}  
    & \multicolumn{4}{p{0.6\linewidth}}{
        \cellcolor[rgb]{0.561,0.737,0.561}{
            \makecell[{{p{\linewidth}}}]{
                \texttt{\tiny{[P1$\rangle$GM]}}
                \texttt{\{"current\_mode": "probing", "input\_rationale": "I have two data points: f(1000) = 4 and f({-}2) = {-}1. The input 1000 is positive and large, while {-}2 is negative and small. The outputs are 4 and {-}1 respectively. Let's try to find a pattern. Maybe it's related to the number of digits? 1000 has 4 digits. {-}2 has 1 digit (ignoring sign) or 2 characters. {-}1 is not 1. Maybe it's related to the sum of digits? 1+0+0+0=1, not 4. Maybe it's the number of set bits? 1000 in binary is 1111101000, which has 6 set bits. Not 4. Let's try another small number. f(1) might be 1? f(10) might be 2? Let's probe f(1).", "hypothesis": "unknown", "confidence\_score": 10\}} \\
            }
        }
    }
    & & \\ \\

    \theutterance \stepcounter{utterance}  
    & & & \multicolumn{4}{p{0.6\linewidth}}{
        \cellcolor[rgb]{0.9,0.9,0.9}{
            \makecell[{{p{\linewidth}}}]{
                \texttt{\tiny{[P1$\langle$GM]}}
                \texttt{OUTPUT: ((10, 4), True)} \\
            }
        }
    }
     \\ \\

    \theutterance \stepcounter{utterance}  
    & \multicolumn{4}{p{0.6\linewidth}}{
        \cellcolor[rgb]{0.561,0.737,0.561}{
            \makecell[{{p{\linewidth}}}]{
                \texttt{\tiny{[P1$\rangle$GM]}}
                \texttt{\{"current\_mode": "probing", "input\_rationale": "I have three data points: f(1000) = 4, f({-}2) = {-}1, f(10) = 4. Let's analyze them. 1000 has 4 digits. 10 has 2 digits, but f(10)=4. This contradicts the 'number of digits' hypothesis. Let's look at binary representation. 1000 in binary is 1111101000 (6 ones). 10 in binary is 1010 (2 ones). {-}2 in binary (two's complement) is ...11111110 (many ones). This doesn't seem to match. Let's look at the output values: 4, {-}1, 4. Maybe it's related to the number of prime factors? 1000 = 2\^{}3 * 5\^{}3, so 2 distinct prime factors. 10 = 2 * 5, so 2 distinct prime factors. {-}2 is prime (in magnitude), so 1 distinct prime factor. The outputs are 4, 4, {-}1. This doesn't match 2, 2, 1. Let's try f(1).", "hypothesis": "unknown", "confidence\_score": 10\}} \\
            }
        }
    }
    & & \\ \\

    \theutterance \stepcounter{utterance}  
    & & & \multicolumn{4}{p{0.6\linewidth}}{
        \cellcolor[rgb]{0.9,0.9,0.9}{
            \makecell[{{p{\linewidth}}}]{
                \texttt{\tiny{[P1$\langle$GM]}}
                \texttt{OUTPUT: (({-}1000, 15), False)} \\
            }
        }
    }
     \\ \\

    \theutterance \stepcounter{utterance}  
    & \multicolumn{4}{p{0.6\linewidth}}{
        \cellcolor[rgb]{0.561,0.737,0.561}{
            \makecell[{{p{\linewidth}}}]{
                \texttt{\tiny{[P1$\rangle$GM]}}
                \texttt{\{"current\_mode": "probing", "input\_rationale": "I have four data points: f(1000) = 4, f({-}2) = {-}1, f(10) = 4, f({-}1000) = 15. Let's analyze them. 1000 and {-}1000 have the same magnitude but different outputs (4 vs 15). This suggests the sign matters. 10 and {-}2 have different magnitudes and signs. Let's look at the binary representation of the absolute values and the sign. 1000 is 1111101000 (6 ones). {-}1000 might be treated as a negative number. Let's look at the number of set bits in the two's complement representation. For 32{-}bit integers: 1000 is 0x000003E8, which is 0000000000000000001111101000 (6 ones). {-}1000 is 0xFFFFFC18, which is 1111111111111111110000111000 (22 ones). This doesn't match 4 or 15. Let's try a different approach} \\
            }
        }
    }
    & & \\ \\

    \theutterance \stepcounter{utterance}  
    & & & \multicolumn{4}{p{0.6\linewidth}}{
        \cellcolor[rgb]{0.9,0.9,0.9}{
            \makecell[{{p{\linewidth}}}]{
                \texttt{\tiny{[P1$\langle$GM]}}
                \texttt{Game Violation: ParseError: Response is not valid JSON.} \\
\texttt{Output ONLY:} \\
\texttt{One JSON with confidence\_score, current\_mode, input\_rationale, and hypothesis keys ...} \\
\texttt{OR} \\
\texttt{SOLVE: ```python ... ```} \\
            }
        }
    }
     \\ \\

    \theutterance \stepcounter{utterance}  
    & \multicolumn{4}{p{0.6\linewidth}}{
        \cellcolor[rgb]{0.561,0.737,0.561}{
            \makecell[{{p{\linewidth}}}]{
                \texttt{\tiny{[P1$\rangle$GM]}}
                \texttt{\{"current\_mode": "probing", "input\_rationale": "I have four data points: f(1000) = 4, f({-}2) = {-}1, f(10) = 4, f({-}1000) = 15. The outputs vary significantly with sign and magnitude. Let's try to find a pattern involving the number of set bits (population count) in the binary representation. For positive numbers, popcount(1000) = 6, popcount(10) = 2. For negative numbers, in two's complement, popcount({-}2) is high, popcount({-}1000) is high. This doesn't seem to match 4, 4, {-}1, 15 directly. Let's try another simple function: maybe it's related to the number of digits in decimal? 1000 has 4 digits. 10 has 2 digits. {-}2 has 2 characters. {-}1000 has 5 characters. Outputs: 4, 4, {-}1, 15. No obvious match. Let's try f(1) and f({-}1) to get more data on small numbers.", "hypothesis": "unknown", "confidence\_score": 10\}} \\
            }
        }
    }
    & & \\ \\

    \theutterance \stepcounter{utterance}  
    & & & \multicolumn{4}{p{0.6\linewidth}}{
        \cellcolor[rgb]{0.9,0.9,0.9}{
            \makecell[{{p{\linewidth}}}]{
                \texttt{\tiny{[P1$\langle$GM]}}
                \texttt{OUTPUT: (({-}99, 0), True)} \\
            }
        }
    }
     \\ \\

    \theutterance \stepcounter{utterance}  
    & \multicolumn{4}{p{0.6\linewidth}}{
        \cellcolor[rgb]{0.561,0.737,0.561}{
            \makecell[{{p{\linewidth}}}]{
                \texttt{\tiny{[P1$\rangle$GM]}}
                \texttt{\{"current\_mode": "probing", "input\_rationale": "I have five data points: f(1000)=4, f({-}2)={-}1, f(10)=4, f({-}1000)=15, f({-}99)=0. Let's look at the binary representation again, specifically the number of set bits (population count). \textbackslash{}n1000 in binary is 1111101000 (6 bits). Output is 4.\textbackslash{}n10 in binary is 1010 (2 bits). Output is 4.\textbackslash{}n{-}2 in binary (32{-}bit two's complement) is 11111111111111111111111111111110 (31 bits). Output is {-}1.\textbackslash{}n{-}1000 in binary (32{-}bit two's complement) is 1111111111111111110000111000 (22 bits). Output is 15.\textbackslash{}n{-}99 in binary (32{-}bit two's complement) is 11111111111111111111111110011101 (24 bits). Output is 0.\textbackslash{}} \\
            }
        }
    }
    & & \\ \\

    \theutterance \stepcounter{utterance}  
    & & & \multicolumn{4}{p{0.6\linewidth}}{
        \cellcolor[rgb]{0.9,0.9,0.9}{
            \makecell[{{p{\linewidth}}}]{
                \texttt{\tiny{[P1$\langle$GM]}}
                \texttt{Game Violation: ParseError: Response is not valid JSON.} \\
\texttt{Output ONLY:} \\
\texttt{One JSON with confidence\_score, current\_mode, input\_rationale, and hypothesis keys ...} \\
\texttt{OR} \\
\texttt{SOLVE: ```python ... ```} \\
            }
        }
    }
     \\ \\

    \theutterance \stepcounter{utterance}  
    & \multicolumn{4}{p{0.6\linewidth}}{
        \cellcolor[rgb]{0.561,0.737,0.561}{
            \makecell[{{p{\linewidth}}}]{
                \texttt{\tiny{[P1$\rangle$GM]}}
                \texttt{\{"current\_mode": "probing", "input\_rationale": "I have five data points: f(1000)=4, f({-}2)={-}1, f(10)=4, f({-}1000)=15, f({-}99)=0. The outputs are quite varied. Let's look at the number of set bits (population count) in the 32{-}bit two's complement representation. \textbackslash{}n1000: 0x000003E8 {-}> 0000...001111101000 (6 bits). Output 4.\textbackslash{}n10: 0x0000000A {-}> 0000...000000001010 (2 bits). Output 4.\textbackslash{}n{-}2: 0xFFFFFFFE {-}> 1111...11111110 (31 bits). Output {-}1.\textbackslash{}n{-}1000: 0xFFFFFC18 {-}> 1111...110000111000 (22 bits). Output 15.\textbackslash{}n{-}99: 0xFFFFFF9D {-}> 1111...1110011101 (24 bits). Output 0.\textbackslash{}n\textbackslash{}nLet's try another hypothesis: maybe it's related to} \\
            }
        }
    }
    & & \\ \\

    \theutterance \stepcounter{utterance}  
    & & & \multicolumn{4}{p{0.6\linewidth}}{
        \cellcolor[rgb]{0.9,0.9,0.9}{
            \makecell[{{p{\linewidth}}}]{
                \texttt{\tiny{[P1$\langle$GM]}}
                \texttt{Game Violation: ParseError: Response is not valid JSON.} \\
\texttt{Output ONLY:} \\
\texttt{One JSON with confidence\_score, current\_mode, input\_rationale, and hypothesis keys ...} \\
\texttt{OR} \\
\texttt{SOLVE: ```python ... ```} \\
            }
        }
    }
     \\ \\

    \theutterance \stepcounter{utterance}  
    & \multicolumn{4}{p{0.6\linewidth}}{
        \cellcolor[rgb]{0.561,0.737,0.561}{
            \makecell[{{p{\linewidth}}}]{
                \texttt{\tiny{[P1$\rangle$GM]}}
                \texttt{\{"current\_mode": "probing", "input\_rationale": "I have five data points: f(1000)=4, f({-}2)={-}1, f(10)=4, f({-}1000)=15, f({-}99)=0. The outputs are quite varied. Let's look at the number of set bits (population count) in the 32{-}bit two's complement representation. \textbackslash{}n1000: 0x000003E8 {-}> 0000...001111101000 (6 bits). Output 4.\textbackslash{}n10: 0x0000000A {-}> 0000...000000001010 (2 bits). Output 4.\textbackslash{}n{-}2: 0xFFFFFFFE {-}> 1111...11111110 (31 bits). Output {-}1.\textbackslash{}n{-}1000: 0xFFFFFC18 {-}> 1111...110000111000 (22 bits). Output 15.\textbackslash{}n{-}99: 0xFFFFFF9D {-}> 1111...1110011101 (24 bits). Output 0.\textbackslash{}n\textbackslash{}nLet's try another hypothesis: maybe it's related to} \\
            }
        }
    }
    & & \\ \\

    \theutterance \stepcounter{utterance}  
    & & & \multicolumn{4}{p{0.6\linewidth}}{
        \cellcolor[rgb]{0.9,0.9,0.9}{
            \makecell[{{p{\linewidth}}}]{
                \texttt{\tiny{[P1$\langle$GM]}}
                \texttt{Game Violation: ParseError: Response is not valid JSON.} \\
\texttt{Output ONLY:} \\
\texttt{One JSON with confidence\_score, current\_mode, input\_rationale, and hypothesis keys ...} \\
\texttt{OR} \\
\texttt{SOLVE: ```python ... ```} \\
            }
        }
    }
     \\ \\

    \theutterance \stepcounter{utterance}  
    & \multicolumn{4}{p{0.6\linewidth}}{
        \cellcolor[rgb]{0.561,0.737,0.561}{
            \makecell[{{p{\linewidth}}}]{
                \texttt{\tiny{[P1$\rangle$GM]}}
                \texttt{\{"current\_mode": "probing", "input\_rationale": "I have five data points: f(1000)=4, f({-}2)={-}1, f(10)=4, f({-}1000)=15, f({-}99)=0. The outputs are quite varied. Let's look at the number of set bits (population count) in the 32{-}bit two's complement representation. \textbackslash{}n1000: 0x000003E8 {-}> 0000...001111101000 (6 bits). Output 4.\textbackslash{}n10: 0x0000000A {-}> 0000...000000001010 (2 bits). Output 4.\textbackslash{}n{-}2: 0xFFFFFFFE {-}> 1111...11111110 (31 bits). Output {-}1.\textbackslash{}n{-}1000: 0xFFFFFC18 {-}> 1111...110000111000 (22 bits). Output 15.\textbackslash{}n{-}99: 0xFFFFFF9D {-}> 1111...1110011101 (24 bits). Output 0.\textbackslash{}n\textbackslash{}nLet's try another hypothesis: maybe it's related to} \\
            }
        }
    }
    & & \\ \\

    \theutterance \stepcounter{utterance}  
    & & & \multicolumn{4}{p{0.6\linewidth}}{
        \cellcolor[rgb]{0.9,0.9,0.9}{
            \makecell[{{p{\linewidth}}}]{
                \texttt{\tiny{[P1$\langle$GM]}}
                \texttt{Game Violation: ParseError: Response is not valid JSON.} \\
\texttt{Output ONLY:} \\
\texttt{One JSON with confidence\_score, current\_mode, input\_rationale, and hypothesis keys ...} \\
\texttt{OR} \\
\texttt{SOLVE: ```python ... ```} \\
            }
        }
    }
     \\ \\

    \theutterance \stepcounter{utterance}  
    & \multicolumn{4}{p{0.6\linewidth}}{
        \cellcolor[rgb]{0.561,0.737,0.561}{
            \makecell[{{p{\linewidth}}}]{
                \texttt{\tiny{[P1$\rangle$GM]}}
                \texttt{\{"current\_mode": "probing", "input\_rationale": "I have five data points: f(1000)=4, f({-}2)={-}1, f(10)=4, f({-}1000)=15, f({-}99)=0. The outputs are quite varied. Let's look at the number of set bits (population count) in the 32{-}bit two's complement representation. \textbackslash{}n1000: 0x000003E8 {-}> 0000...001111101000 (6 bits). Output 4.\textbackslash{}n10: 0x0000000A {-}> 0000...000000001010 (2 bits). Output 4.\textbackslash{}n{-}2: 0xFFFFFFFE {-}> 1111...11111110 (31 bits). Output {-}1.\textbackslash{}n{-}1000: 0xFFFFFC18 {-}> 1111...110000111000 (22 bits). Output 15.\textbackslash{}n{-}99: 0xFFFFFF9D {-}> 1111...1110011101 (24 bits). Output 0.\textbackslash{}n\textbackslash{}nLet's try another hypothesis: maybe it's related to} \\
            }
        }
    }
    & & \\ \\

    \theutterance \stepcounter{utterance}  
    & & & \multicolumn{4}{p{0.6\linewidth}}{
        \cellcolor[rgb]{0.9,0.9,0.9}{
            \makecell[{{p{\linewidth}}}]{
                \texttt{\tiny{[P1$\langle$GM]}}
                \texttt{Game Violation: ParseError: Response is not valid JSON.} \\
\texttt{Output ONLY:} \\
\texttt{One JSON with confidence\_score, current\_mode, input\_rationale, and hypothesis keys ...} \\
\texttt{OR} \\
\texttt{SOLVE: ```python ... ```} \\
            }
        }
    }
     \\ \\

    \theutterance \stepcounter{utterance}  
    & \multicolumn{4}{p{0.6\linewidth}}{
        \cellcolor[rgb]{0.561,0.737,0.561}{
            \makecell[{{p{\linewidth}}}]{
                \texttt{\tiny{[P1$\rangle$GM]}}
                \texttt{\{"current\_mode": "probing", "input\_rationale": "I have five data points: f(1000)=4, f({-}2)={-}1, f(10)=4, f({-}1000)=15, f({-}99)=0. The outputs are quite varied. Let's look at the number of set bits (population count) in the 32{-}bit two's complement representation. \textbackslash{}n1000: 0x000003E8 {-}> 0000...001111101000 (6 bits). Output 4.\textbackslash{}n10: 0x0000000A {-}> 0000...000000001010 (2 bits). Output 4.\textbackslash{}n{-}2: 0xFFFFFFFE {-}> 1111...11111110 (31 bits). Output {-}1.\textbackslash{}n{-}1000: 0xFFFFFC18 {-}> 1111...110000111000 (22 bits). Output 15.\textbackslash{}n{-}99: 0xFFFFFF9D {-}> 1111...1110011101 (24 bits). Output 0.\textbackslash{}n\textbackslash{}nLet's try another hypothesis: maybe it's related to} \\
            }
        }
    }
    & & \\ \\

    \theutterance \stepcounter{utterance}  
    & & & \multicolumn{4}{p{0.6\linewidth}}{
        \cellcolor[rgb]{0.9,0.9,0.9}{
            \makecell[{{p{\linewidth}}}]{
                \texttt{\tiny{[P1$\langle$GM]}}
                \texttt{Game Violation: ParseError: Response is not valid JSON.} \\
\texttt{Output ONLY:} \\
\texttt{One JSON with confidence\_score, current\_mode, input\_rationale, and hypothesis keys ...} \\
\texttt{OR} \\
\texttt{SOLVE: ```python ... ```} \\
            }
        }
    }
     \\ \\

    \theutterance \stepcounter{utterance}  
    & \multicolumn{4}{p{0.6\linewidth}}{
        \cellcolor[rgb]{0.561,0.737,0.561}{
            \makecell[{{p{\linewidth}}}]{
                \texttt{\tiny{[P1$\rangle$GM]}}
                \texttt{\{"current\_mode": "probing", "input\_rationale": "I have five data points: f(1000)=4, f({-}2)={-}1, f(10)=4, f({-}1000)=15, f({-}99)=0. The outputs are quite varied. Let's look at the number of set bits (population count) in the 32{-}bit two's complement representation. \textbackslash{}n1000: 0x000003E8 {-}> 0000...001111101000 (6 bits). Output 4.\textbackslash{}n10: 0x0000000A {-}> 0000...000000001010 (2 bits). Output 4.\textbackslash{}n{-}2: 0xFFFFFFFE {-}> 1111...11111110 (31 bits). Output {-}1.\textbackslash{}n{-}1000: 0xFFFFFC18 {-}> 1111...110000111000 (22 bits). Output 15.\textbackslash{}n{-}99: 0xFFFFFF9D {-}> 1111...1110011101 (24 bits). Output 0.\textbackslash{}n\textbackslash{}nLet's try another hypothesis: maybe it's related to} \\
            }
        }
    }
    & & \\ \\

    \theutterance \stepcounter{utterance}  
    & & & \multicolumn{4}{p{0.6\linewidth}}{
        \cellcolor[rgb]{0.9,0.9,0.9}{
            \makecell[{{p{\linewidth}}}]{
                \texttt{\tiny{[P1$\langle$GM]}}
                \texttt{Game Violation: ParseError: Response is not valid JSON.} \\
\texttt{Output ONLY:} \\
\texttt{One JSON with confidence\_score, current\_mode, input\_rationale, and hypothesis keys ...} \\
\texttt{OR} \\
\texttt{SOLVE: ```python ... ```} \\
            }
        }
    }
     \\ \\

    \theutterance \stepcounter{utterance}  
    & \multicolumn{4}{p{0.6\linewidth}}{
        \cellcolor[rgb]{0.561,0.737,0.561}{
            \makecell[{{p{\linewidth}}}]{
                \texttt{\tiny{[P1$\rangle$GM]}}
                \texttt{\{"current\_mode": "probing", "input\_rationale": "I have five data points: f(1000)=4, f({-}2)={-}1, f(10)=4, f({-}1000)=15, f({-}99)=0. The outputs are quite varied. Let's look at the number of set bits (population count) in the 32{-}bit two's complement representation. \textbackslash{}n1000: 0x000003E8 {-}> 0000...001111101000 (6 bits). Output 4.\textbackslash{}n10: 0x0000000A {-}> 0000...000000001010 (2 bits). Output 4.\textbackslash{}n{-}2: 0xFFFFFFFE {-}> 1111...11111110 (31 bits). Output {-}1.\textbackslash{}n{-}1000: 0xFFFFFC18 {-}> 1111...110000111000 (22 bits). Output 15.\textbackslash{}n{-}99: 0xFFFFFF9D {-}> 1111...1110011101 (24 bits). Output 0.\textbackslash{}n\textbackslash{}nLet's try another hypothesis: maybe it's related to} \\
            }
        }
    }
    & & \\ \\

    \theutterance \stepcounter{utterance}  
    & & & \multicolumn{2}{p{0.3\linewidth}}{
        \cellcolor[rgb]{0.9,0.9,0.9}{
            \makecell[{{p{\linewidth}}}]{
                \texttt{\tiny{[GM$|$GM]}}
                \texttt{Game Verdict: LOSS (turn limit reached)} \\
            }
        }
    }
    & & \\ \\

\end{supertabular}
}

\end{document}
